\section{Implementation}
\label{sec:impl}

\systemname\ is implemented in Python and PyTorch~\cite{paszke2019pytorch} and integrates with Hugging Face Transformers~\cite{wolf2020transformers}. The implementation is divided into a model-independent control plane and thin model adapters. This separation keeps scheduling and memory-safety logic independent of a particular router or expert MLP layout.

\subsection{Control Plane}

\paragraph{Observation and scheduling.}
Model adapters report the router's top-$k$ expert identifiers and routing weights to a \texttt{RouterObserver}. A \texttt{HotnessTracker} accumulates per-layer counters and updates the EMA state. At a scheduling tick, \texttt{PrecisionScheduler} applies the per-layer capacity, hysteresis margin, and transition-rate limit to produce typed promotion and demotion requests. The scheduler never mutates a live expert handle.

The calibration entry point consumes at least 128 uniquely identified prompts from train, validation, development, or explicitly designated calibration splits; test splits are refused. It forces the all-low representation, disables online reassignment, and accumulates each prompt's token-normalized routed probability mass in a separate FP64 counter. Dividing by prompt count yields an order-invariant calibration score, while the FP32 EMA remains dedicated to online adaptation. The entry point emits a full stable ranking per layer and hashes both the source file and selected identifiers. Formal runs require this schema-v2 ranking artifact to match the exact checkpoint and model contract. The runtime records the realized prefix set, allowing the audit to prove that full, static, blocking, and sensitivity runs began from the intended common map.

\paragraph{Handles and transition lifecycle.}
\texttt{ExpertRegistry} stores one versioned handle per $(\ell,e)$ pair. A handle records its precision tier, packed quantization metadata, originating pool block, byte footprint, reservation, active-reader count, and latest event per compute stream. \texttt{TransitionEngine} uses a bounded worker pool to execute four stages: obtain the prepared host representation, reserve and allocate the destination block, copy the packed bytes, and publish the new handle. A failed reservation or exhausted pool leaves the previous handle active. After publication, the engine waits for all leases on the old object to drain and then for its per-stream last-use events before returning the block to its originating resident or global staging pool. If an adapter cannot provide a precise event, the implementation conservatively falls back to device synchronization; both paths are counted, and formal asynchronous artifacts require zero global-sync fallbacks.

\paragraph{Memory accounting.}
\texttt{BudgetTracker} maintains committed and pending bytes separately for the two precision tiers and enforces a cross-tier live-byte cap. Its \texttt{try\_reserve}, \texttt{commit}, and \texttt{release} operations are serialized by a lock, and a staging cap bounds pending work. The allocator pre-allocates disjoint per-layer, per-tier resident pools plus a global transient pool; transitions invoke no global device allocation. A compare-and-swap registry update repatriates staging-backed handles when the paired swap creates a resident hole. Because expert shapes may differ across models, block sizes and counts are derived from the packed representation rather than nominal bit width.

\subsection{Quantized Representations and Dispatch}

\paragraph{Packed representation.}
The quantization layer exposes a common \texttt{PackedTensor} abstraction for FP16, group-wise symmetric INT4 (128-weight groups), and group-wise symmetric INT2 (64-weight groups). It stores packed bytes, FP16 scales, tensor dimensions, group size, and the physical storage charged by admission control. Reference pack, unpack, and dequantize-then-GEMM paths support correctness testing. A compatibility backend invokes AutoRound's per-tensor symmetric quantizer and converts its output to the same representation; this path implements symmetric rounding, not AutoRound's activation-calibrated optimization loop.

\paragraph{Device kernels and workspace.}
INT4 dispatch uses PyTorch's fused int4pack matrix multiplication. INT2 dispatch uses a Triton kernel that unpacks values within the matrix multiplication and therefore avoids a full FP16 expert temporary. During an INT4 transition, the kernel-native layout replaces the canonical transfer layout within the same pool block, and its scale-and-zero storage remains charged to the handle. Conversion briefly creates a nibble-swapped input, a native-layout output, and scale/zero tensors. Startup consequently reserves a conservative two-expert INT4 conversion workspace per admitted transition before solving the resident allocation. Admission accounts for these physical layouts and workspaces, rather than nominal weight bits, and each artifact records the selected backend. A formal run aborts if the requested format lacks a device kernel that satisfies this envelope.

\paragraph{Model adapters.}
The Qwen3-MoE, Qwen3-Next, and Phi-3.5-MoE adapters acquire a handle immediately before routed-expert dispatch and release its lease after the final routed kernel. Gate, up, and down projections retain separate packed metadata but participate in one expert transition. Qwen3-MoE switches from Transformers' grouped-expert shortcut to handle-aware dispatch, because the shortcut bypasses per-expert resolution. Qwen3-Next's shared expert remains on the fixed dense path and is charged to the dense reservation. Once attached, an adapter fails closed on a missing or incomplete handle; native dispatch remains available only before attachment and in equivalence tests.

\paragraph{Startup path.}
Before a dynamic experiment, the runtime verifies that the declared layer count, routed-expert count, and top-$k$ match the checkpoint. It loads the source on CPU, packs both tiers in 16-expert chunks, releases native expert tensors layer by layer, and returns freed host arenas to the operating system. Only the remaining dense model is then moved to one specified GPU. Startup measures all remaining parameter and buffer bytes, raises an underestimated dense reservation, sizes conversion workspace from the chosen tier pair, and allocates the fixed pools. The run is rejected unless one device can hold the pools together with the declared KV-cache, activation, workspace, and safety reserves. Integration validation additionally requires that the attached model both emit routing observations and consume registered handles.

\subsection{Testing and Failure Semantics}

\paragraph{Correctness tests.}
CPU tests cover bit packing, quantization error, byte accounting, paired rate limiting, reservation races, global staging allocation, publish-after-copy ordering, lease-and-event-fenced reclamation, and repeated transitions without block leakage. Integration tests compare the Phi-MoE, Qwen3-MoE, and upstream Qwen3-Next adapters with native forward execution when handles reference equivalent weights. The Qwen3-Next test also verifies that releasing routed source tensors preserves its shared expert.

\paragraph{Experiment failure semantics.}
The evaluation CLI distinguishes incorrect model outputs from infrastructure failures and applies one named paper protocol. Paired comparisons require identical dataset fingerprints and sample identifiers before computing exact McNemar tests with Holm correction. The MoE-Infinity adapter binds the official repository commit, recursive source-tree manifest, and model snapshot, and accepts a run only when both offloaded expert state and measured prefetch activity are present. Transition artifacts record rejection causes, bytes copied, stage timings, pool occupancy, workspace reservations, and final budget counters.

During each formal timed generation, a 2\,ms NVML poller maps the selected CUDA device to its physical UUID and records current-process HBM use, including native allocations outside PyTorch's counters. Idle foreign HBM residency is recorded separately, while any observed foreign-process utilization invalidates the sample. Finally, a registrar hashes each artifact with its reproduction command. The manuscript audit enumerates every empirical and statistical claim and rejects missing, modified, dirty-tree, method-inconsistent, incompletely sampled, or memory-unaccounted evidence. Optional backends and large-model tests remain outside the default suite, so an unavailable quantizer or checkpoint cannot silently produce a passing experiment.
